% Retrospective evidence qualification table (Task6-R), D2T-RNA v7.
% This is the canonical table included in manuscript Sec \ref{sec:retro}.
\caption{Retrospective evidence qualification. \textbf{Question:} how do the registered datasets
close, and what may/may not be reported? \textbf{Result:} add, SAM-III, and RORC close fail-closed
with distinct terminal outcomes. \textbf{Interpretation:} data are used only as fixed realized catalogs
for a fail-closed audit. \textbf{Boundary:} no categorical replay, no quantitative T2 validation, no
future cost saving, no independent validation.}
\label{tab:retro}
\begin{tabular}{@{}p{2.0cm}p{3.2cm}p{3.4cm}p{3.2cm}p{3.2cm}@{}}
\toprule
Dataset & Registered role & R2 outcome & What can be reported & What cannot be claimed \\
\midrule
add/RMDB & retrospective full-matrix compression if qualified &
  NOT\_COMPARABLE\_BY\_REGISTERED\_OBSERVATION\_MODEL &
  modality limitation; raw provenance; fail-closed audit &
  categorical replay; quantitative T2 validation; future cost saving \\
\addlinespace
SAM-III/GSE278422 & modality-transfer diagnostic &
  NOT\_COMPARABLE\_BY\_REGISTERED\_ACTION\_SPACE &
  action/modality incompatibility; diagnostic failure &
  unified benchmark; quantitative transfer \\
\addlinespace
RORC & exposed misspecification stress &
  NOT\_APPLICABLE &
  explicit terminal audit and boundary semantics &
  independent validation; unseen OOD; third-state discovery \\
\bottomrule
\end{tabular}